Shelfmark 149, scan 135.

\scanpage{135}
\mspage{67}

We shall therefore assume that (37) does indeed define a structure of \unsure{extensions|extension}, and we shall have a homomorphism of exact sequences

\deleted{\unsure{(38)|(37)}}
\[
\tag{38}
\diagram{
\begin{tikzcd}[row sep=large,column sep=large]
1 \arrow[r,dash] & T^J \arrow[r] \arrow[d] & \mathcal E_{U_{J,\alpha}}^{\bar L} \arrow[r] \arrow[d] & \mathcal E_{D_{J,\alpha}}^{\tilde D_{J,\alpha}} \arrow[r] \arrow[d] & 1 \\
1 \arrow[r,dash] & \pi_{U/K}^{\bar L} \arrow[r] & \mathcal E_U^{\bar L} \arrow[r] & \mathcal E_K^{\bar K} \arrow[r] & 1
\end{tikzcd}
}
\]
--- where we have corrected our aim, bearing in mind that (37) only makes sense as it stands if $D_J$ is connected, \inserted{that is, $X_J$ connected} that otherwise we must look at the connected components $X_{J,\alpha}$ of $X_J$, inducing the \unsure{connected|connected} components $D_{J,\alpha}$ of $D_J$, and $U_{J,\alpha}$ of $U_J$, and write the exact sequence for each component. Of course, there are more intrinsic things: write with the \deleted{\unsure{Galois|geometric} étale}\inserted{multigaloisian} topoi, and their fundamental groupoids....\deleted{\unsure{classes of} $\mathcal E_U^{\bar L}$ \unsure{conjugacy|conjugates}}

We thus see that the homomorphisms
\[
\tag{39}
T^J\longrightarrow\unsure{\pi|\mathcal E}_{U/K}^{\bar L}
\]
induced by the different $x\in D_J$ --- \unsure{each} being in fact, when we return to the precise definition, only \inserted{defined} modulo operations \unsure{induced|direct} by inner automorphisms of $\mathcal E_U^{\bar L}$, induced by

\margin{134}
